\section{Semilinear Clifford lifts and the even-dimensional Spin boundary}\label{semilinear-clifford-lifts-and-the-even-dimensional-spin-boundary}

\subsection{Status and exact scope}\label{status-and-exact-scope}

New author derivation, D-OL19--21. Independent and root whole-proof reviews are pending. The field has odd characteristic. The lifting lemma is valid in either parity in dimension at least 22; the final support and source theorems concern even natural dimension at least 1136. No source exclusion in odd natural dimension, characteristic two, or small-rank triality is made here.

We use the complete natural-group arguments D-OL06--09 and D-OL12--18, retaining their exact orbit estimates, finite-field descent, and actual-source hypotheses. The new points are: a semilinear lift acting on every represented central quotient of Spin; unique lifts of the odd-order radicals; and a degree estimate that applies to every nontrivial Spin module. The proof does not assume that a module descends to the natural Omega group or extends to a full orthogonal group.

Write

\[
 \beta=\frac{4997}{30000},\qquad
 \delta=\frac{1-(19/20)^9}{2},\qquad
 \nu(T)=a-\max_\lambda\dim\ker(T-\lambda I)
 \tag{SP.1}
\]

for an operator on a complex vector space of dimension \(a\).

\subsubsection{Exact published group inputs}\label{exact-published-group-inputs}

Let \(F=\mathbf F_q\), \(q\) odd, and let \((V,Q)\) be nondegenerate. The finite Spin group has its usual Clifford realization

\[
 S=\operatorname{Spin}(V)
 =\{g\in C^0(V,Q)^*:gVg^{-1}=V,\ \tau(g)g=1\},
 \qquad
 1\longrightarrow\{\pm1\}\longrightarrow S
 \mathop{\longrightarrow}^{\varphi}\Omega(V)\longrightarrow1,
 \tag{SP.2}
\]

where \(\tau\) reverses Clifford words. In the dimensions used here \(S\) is perfect and quasisimple. Precise published locators are Larsen--Shalev--Tiep, \emph{The Waring problem for finite simple groups}, Annals 174 (2011), proof of Theorem 4.3.6, Step 1, pp.1918--1919: the second exact sequence on p.1919 is (SP.2), and the following paragraph explicitly identifies \(K_i=\operatorname{Spin}(W_i)\) as quasisimple for \(\dim W_i\ge5\). \cite[proof of Theorem~4.3.6, Step~1, pp.~1918--1919]{HOS:LST2011} This is an invoked published structural input, not an inference that an arbitrary commutator subgroup is perfect. The Clifford definition and exact sequence also appear in Liebeck--O'Brien--Shalev--Tiep, \emph{Commutators in finite quasisimple groups}, Bulletin LMS 43 (2011), Section 4; the freshly read author version has that passage on pp.10--11 \cite[Section~4, author pp.~10--11]{SF:LOST2011}.

The simple quotient is \(\overline S=P\Omega(V)\). One can identify its kernel without using a table of Spin centers. If \(g\in S\) maps into \(Z(\Omega(V))\), then \([g,s]\in\{\pm1\}\) for every \(s\in S\). The map \(s\mapsto[g,s]\) is a homomorphism because its values are central, and perfectness makes it trivial. Thus

\[
 Z(S)=\varphi^{-1}(Z(\Omega(V))),\qquad
 S/Z(S)\simeq P\Omega(V).
 \tag{SP.3}
\]

The complete natural semilinear lifting argument D-OL07 applies to this simple quotient in dimension at least 22. It uses the published outer-order table of Li--Xia, Section 2.2.2, Table 2.1, \cite[Section~2.2.2, Table~2.1]{HES:LiXia2022}, and explicitly proves that its full outer image is realized by natural semilinear similitudes, with scalar kernel. We retain that input and its triality restriction.

\subsection{D-OL19. Every compatible represented Spin normalizer has a natural lift}\label{d-ol19.-every-compatible-represented-spin-normalizer-has-a-natural-lift}

\textbf{Lemma.} Suppose \(q\) is odd and \(d=\dim V\ge22\), with either parity. Let \(\rho:S\to\mathrm{GL}(M)\) be a nontrivial irreducible complex representation, with arbitrary central kernel. Every operator normalizing \(\rho(S)\) is compatible with a natural semilinear similitude of \((V,Q)\). The induced projective orders agree. This assertion includes modules that do not descend through \(\varphi\).

\textbf{Proof.} We first give the natural lift explicitly. Let \(H\) be \(\sigma\)-semilinear and have multiplier \(\mu\in F^*\), so

\[
 Q(Hv)=\mu\,\sigma(Q(v)).
 \tag{SP.4}
\]

On even tensor words define a \(\sigma\)-semilinear map by

\[
 \Phi_H(v_1\cdots v_{2k})
 =\mu^{-k}Hv_1\cdots Hv_{2k}.
 \tag{SP.5}
\]

Here a Clifford word has no well-defined integer degree, so well-definedness must be checked before using this formula. In the tensor algebra, an even word containing a relation \(v\otimes v-Q(v)\) maps to a scalar multiple of the corresponding relation \(Hv\otimes Hv-Q(Hv)\). Indeed the two word lengths differ by two, and the extra factor \(\mu^{-1}\) cancels \(Q(Hv)=\mu\sigma Q(v)\). Thus the defining ideal in the even tensor algebra is preserved. The inverse is the construction for \(H^{-1}\). Products of even words acquire exactly the product of their scaling factors, so (SP.5) is a semilinear algebra automorphism of \(C^0(V,Q)\).

This is also the construction and complete ideal calculation in Wonenburger, \emph{The Clifford Algebra and the Group of Similitudes}, Canadian J. Math. 14 (1962), Part II, pp.48--49. \cite[Part~II, pp.~48--49]{SF:Wonenburger1962} Its use here requires no extension to the entire Clifford algebra; Wonenburger's separate Theorem 3 imposes conditions on such an extension, and those conditions are not being assumed.

For the needed vector covariance, define on odd words a semilinear map

\[
 \Psi_H(v_1\cdots v_{2k+1})
 =\mu^{-k}Hv_1\cdots Hv_{2k+1}.
 \tag{SP.6}
\]

The same two-length calculation preserves the odd part of the defining ideal. For \(a,b\in C^0\) and \(z\in C^1\), counting lengths and then passing to the quotient gives

\[
 \Psi_H(azb)=\Phi_H(a)\Psi_H(z)\Phi_H(b),
 \qquad \Psi_H|_V=H.
 \tag{SP.7}
\]

The pair \((\Phi_H,\Psi_H)\) is an algebra automorphism on the even part and a compatible map on its odd bimodule. It is not claimed to be an automorphism of the whole Clifford algebra: the product of two odd images has the additional factor \(\mu\).

Word reversal commutes with \(\Phi_H\), because reversing a word changes neither its length nor its scaling factor. If \(g\in S\), then

\[
 \tau(\Phi_H(g))\Phi_H(g)=1,\qquad
 \Phi_H(g)H(v)\Phi_H(g)^{-1}=H(gvg^{-1}).
 \tag{SP.8}
\]

Thus \(\Phi_H(g)\) is again in \(S\), and the inverse construction proves that \(\alpha_H=\Phi_H|_S\) is a group automorphism. Its natural image satisfies

\[
 \varphi(\alpha_H(g))=H\varphi(g)H^{-1}.
 \tag{SP.9}
\]

If \(H_1,H_2\) have semilinear parts \(\sigma_1,\sigma_2\), their composite multiplier is \(\mu_1\sigma_1(\mu_2)\). Substitution in (SP.5) gives \(\alpha_{H_1H_2}=\alpha_{H_1}\alpha_{H_2}\). Replacing \(H\) by a nonzero linear scalar multiple changes \(\mu\) by the scalar's square and leaves (SP.5) unchanged. Hence this lift depends only on the projective semilinear similitude.

Put \(R=\rho(S)\). Quasisimplicity shows that \(\ker\rho\subset Z(S)\). Moreover \(Z(R)=\rho(Z(S))\): if a lift \(s\in S\) has central image in \(R\), all commutators \([s,g]\) lie in the central kernel and again form a homomorphism from the perfect group \(S\); it is trivial. Therefore \(R/Z(R)\simeq\overline S\).

Let \(T\) normalize \(R\), and write \(\gamma\) for its conjugation automorphism of \(R\). Choose a natural semilinear similitude \(H\) realizing the induced automorphism of \(\overline S\), using D-OL07. The two homomorphisms \(\gamma\circ\rho\) and \(\rho\circ\alpha_H\) have the same projection to \(\overline S\). Consequently

\[
 c(s)=\gamma(\rho(s))\,\rho(\alpha_H(s))^{-1}\in Z(R)
 \tag{SP.10}
\]

is a homomorphism \(S\to Z(R)\). For example, the multiplicativity follows by moving the central factor \(c(t)\) past the two images of \(s\). Perfectness implies \(c=1\). We obtain the exact compatibility

\[
 T\rho(s)T^{-1}=\rho(\alpha_H(s))\quad(s\in S).
 \tag{SP.11}
\]

In particular \(\alpha_H\) preserves \(\ker\rho\). Preservation of a chosen central kernel has been proved from the actual represented automorphism; it was not assumed and no classification of the different Spin centers was needed.

Finally, an automorphism of the perfect group \(R\) inducing the identity on \(R/Z(R)\) differs from the identity by a central homomorphism and is the identity. Schur's lemma then says that its implementing operator is scalar. On the natural side D-OL07 proves that a semilinear similitude inducing the identity on \(P\Omega(V)\) is a linear scalar. Applying these two observations to every power proves equality of the projective orders of \(T\) and \(H\). This proves the lemma. \(\square\)

\subsection{D-OL20. Radicals and all the support estimates transfer to Spin}\label{d-ol20.-radicals-and-all-the-support-estimates-transfer-to-spin}

\textbf{Lemma.} Let \(q\) be odd, \(d\ge1136\) even, and let \(M\) be any nontrivial irreducible complex Spin module of degree \(a\). A finite-projective-order nonscalar compatible normalizer satisfies

\[
 \nu(T)\ge\beta a.
 \tag{SP.12}
\]

If \(q\ge5\), it satisfies \(\nu(T)>\delta a\). If \(q=3\) and \(\nu(T)\le\delta a\), its natural projective lift is a reflection and has projective order two. Also

\[
 a\ge q^{d-3}-q^{\lceil(d-2)/2\rceil}-1>360.
 \tag{SP.13}
\]

\textbf{Proof: unique lifts, generation, and multiplicities.} Let \(U\) be any one- or two-plane unipotent radical used in D-OL06 or D-OL16. It is a group of odd prime-power order. The preimage \(\varphi^{-1}(U)\) is a central extension with kernel \(\{\pm1\}\). A Sylow subgroup for that odd prime maps isomorphically onto \(U\). Denote it by \(\widetilde U\). Every element of the preimage is a central element times an element of \(\widetilde U\), so \(\widetilde U\) is normal and is the unique such Sylow subgroup. This proves existence, uniqueness and functoriality of the lift without a splitting choice.

If \(U,U'\) are opposite radicals that generate \(\Omega(V)\), let \(L=\langle\widetilde U,\widetilde U'\rangle\). Its natural image is all of \(\Omega(V)\), so \(S=\{\pm1\}L\). The central kernel makes \(L\) normal and the quotient abelian. Perfectness of \(S\) gives \(L=S\). Natural conjugacies between these radicals lift through \(\varphi\) and therefore conjugate their unique lifts. On a nontrivial irreducible Spin module the two fixed spaces have equal dimension and zero intersection, giving the same bound \(t_0\le a/2\).

For a natural Levi subgroup \(\Omega(W)\), use its full inverse image under \(\varphi\). This suffices for all weight-orbit arguments: every natural Levi element has a lift, and conjugation on \(\widetilde U\) is exactly its prescribed action on \(U\). Weight multiplicities are constant on each orbit. No assertion that the represented module factors through this Levi quotient is made. Likewise (SP.9) and uniqueness imply that \(\alpha_H\) acts on each invariant lifted radical exactly by the natural matrix formulas. There is no additional central character or phase from the double cover, because the odd-order radical maps isomorphically onto its image.

It follows in particular that the full two-plane carrier proof D-OL16 applies to every Spin module: the same Heisenberg group, central character sectors, Weyl basis, fixed quotient dimension, and tensor multiplicity operator occur. For a natural nonsquare projective involution its bound is \(\nu(T)\ge7a/32\).

\textbf{Proof: the normalizer cases.} We give the transfer for every case, so that the ordinary boundary is removed only after checking each use of it. Replace \(T\) by a nonscalar prime-order projective power as in D-OL17; this cannot increase \(\nu\). By D-OL19 its natural lift has the same projective order.

If the field part is nontrivial, it has that prime order. The arbitrary-field norm and form descent in D-OL09 supplies the invariant hyperbolic pair and Frobenius action used in its line-radical count. The unique lift preserves exactly this action, the full inverse-image Levi supplies all orbits, and the bound is \(\nu(T)\ge3a/16>\delta a\).

For a linear odd-prime projective order \(\ell\), normalize the natural lift exactly as before. From \(H^\ell=cI\), with multiplier \(\mu\), one has \(\mu^\ell=c^2\). Set \(z=c/\mu^{(\ell-1)/2}\). Then \(z^2=\mu\), \(z^\ell=c\), and \(g=H/z\) is an odd-order isometry. Its determinant and spinor norm are trivial, so \(g\in\Omega(V)\). Choose \(\widetilde g\in S\) above it. Because \(g\) is an isometry, \(v\mapsto gv\) extends to an automorphism of the whole Clifford algebra whose even restriction is \(\Phi_g\). Conjugation by \(\widetilde g\) agrees with that automorphism on the generating vector space and hence everywhere. Scalar invariance of (SP.5) gives \(\alpha_H=\operatorname{Int}(\widetilde g)\), so \(T\) is a scalar multiple of \(\rho(\widetilde g)\).

When \(g\) fixes a hyperbolic pair, the complementary nontrivial natural action has rank at least two; the all-rank line-radical estimate yields \(7a/32\). When no such pair is fixed, the exact fixed-space argument in D-OL17 gives restricted support against \(\{1,-1\}\) at least \((d-2)/2\ge567\). LST Theorem 4.3.6 and Remark 4.3.7 apply directly to the irreducible Spin character of \(M\). The chosen lift may have an extra central factor; this changes neither its natural action nor the absolute character bound. Thus

\[
 \frac{|\operatorname{tr}\rho(\widetilde g)|}{a}
 <(19/20)^9,
 \qquad \nu(T)>\delta a.
 \tag{SP.14}
\]

For a linear projective involution with \(H^2=cI\), the nonsquare case was covered by the transferred two-plane proof. If \(c\) is square, normalize \(H^2=I\). For multiplier one, the eigenspaces are nondegenerate and one contains a hyperbolic pair. The same rank-at-least-two count applies unless the projective natural action is a reflection. In that last case the exact rank-one fraction from D-OL17 is less than \(1/4\) for \(q\ge5\), or less than \(1-4\beta\) for \(q=3\), giving respectively \(3a/16\) or \(\beta a\). For multiplier minus one, the opposite maximal singular eigenspaces give the same line-radical action and polar-dual action as in D-OL17. The fixed weights have norm zero and their orbit fraction is less than \(1/8\), giving \(7a/32\). Every action just used is the action on the unique lifted radical, so no ordinary-module hypothesis remains.

This proves the universal bound and the stronger bound for \(q\ge5\). Over \(F_3\), suppose \(\nu(T)\le\delta a\). The projective order has no odd prime factor, because a corresponding prime-order power would contradict (SP.14). Its involution power is a natural projective reflection. If its projective order were \(2^s\) with \(s\ge2\), the natural lift \(H\) would preserve the reflection's nonsingular line. If \(H\) acts there by \(\lambda\), its multiplier is \(\lambda^2\), so \(H/\lambda\) is an isometry. Its even power is a scalar isometry projectively equal to a reflection. In even dimension such an isometry is \(r\) or \(-r\), both of determinant minus one, whereas an even power has determinant one. This contradiction proves the order-two reflection conclusion. The parity restriction is essential at this step and is retained.

\textbf{Proof: the degree bound.} Use one full line radical \(U\), whose natural parameter space has dimension \(N=d-2\). If every weight of \(\widetilde U\) on \(M\) were trivial, its opposite conjugate would also act trivially. Their generation of \(S\) would make \(M\) trivial. Thus a nonzero weight occurs. The full inverse image of the Levi \(\Omega(W)\) moves it through its entire natural nonzero vector orbit. The exact quadratic-vector count and transitivity used in D-OL06 give an orbit of size at least \(q^{N-1}-q^{\lceil N/2\rceil}-1\). Every such weight occurs with positive multiplicity, proving (SP.13). This argument applies to faithful Spin modules and uses no ordinary GLT degree theorem. \(\square\)

\subsection{D-OL21. The even-dimensional odd-characteristic Spin source is excluded}\label{d-ol21.-the-even-dimensional-odd-characteristic-spin-source-is-excluded}

\textbf{Theorem.} In the actual inducing-representation and base-cover setting of D-OL18, allow each normal factor to be any represented central quotient of an odd-characteristic Spin group of even natural dimension at least 1136, and allow every nontrivial irreducible complex module of that factor. Then there is no rank-nine sixfold source in this case. Ordinary descent to the natural Omega factors is not required.

\textbf{Proof.} Let \(e=\prod_i a_i\) be the inducing degree. The new degree estimate gives \(a_i>360\) and hence \(e>360\). The same spectral convolution and factor-permutation arguments of D-OL18 apply to these actual modules. Scalar inertia operators different from one move the entire module. Call a nonidentity actual inertia good if its moving rank is at least \(\delta e\), and bad otherwise.

A bad inertia is factor preserving. Each active factor has field \(F_3\), and its natural projective lift is a reflection of order two by D-OL20. On a normalized active factor the two eigenspace proportions lie in \([\beta,1-\beta]\). Two active factors already have projective support at least

\[
 \bigl(1-(1-\beta)^2-\beta^2\bigr)e
 =2\beta(1-\beta)e>\tfrac14e>\delta e.
 \tag{SP.15}
\]

Thus a bad inertia has exactly one active factor. Its actual scalar must make an eigenvalue equal to one; otherwise it would move all \(e\) dimensions. In every case its actual moving rank is at least \(\beta e\).

The unchanged actual Euler identity is

\[
 6=2e(h-1)+D.
 \tag{SP.16}
\]

Because \(e>360\), it gives \(h\le1\). For \(h=1\), \(D=6\) allows at most six marked points and hence period rank at most six. A rank-nine case must therefore have \(h=0\), at least twelve marked points, and \(D=2e+6\). The exact constants satisfy

\[
 13\beta-2>\tfrac17,\qquad
 11\beta+\delta-2>\tfrac1{60}.
 \tag{SP.17}
\]

Thirteen marked points would imply \(e<42\), while twelve with even one good inertia would imply \(e<360\). Consequently the only putative case has exactly twelve bad inertias.

Those actual inertia elements generate the actual inducing group on the sphere. All are factor preserving. Fix a normal factor. Its induced projection contains the full simple inner group \(P\Omega(V_i)\); this follows from the containment of the represented factor in the inducing group and does not require a choice of central splitting. If \(q_i\ne3\), all twelve bad inertias project trivially to it, an immediate contradiction. If \(q_i=3\), their projections are identities or natural projective reflections. Choose at most twelve natural reflections representing them, with nonsingular normals \(v_j\). Their common fixed space

\[
 U=\bigcap_j v_j^\perp,
 \qquad \dim U\ge d_i-12>0
 \tag{SP.18}
\]

is fixed pointwise by every word in those reflections. Every element of the natural group \(\Omega(V_i)\) is projectively such a word, so it acts by an original-field scalar on \(U\). These scalars form a homomorphism \(\Omega(V_i)\to F_3^*\); perfectness makes it trivial. This would make \(U\) a nonzero fixed space of \(\Omega(V_i)\), contradicting its opposite full line radicals and their generation. This argument uses the natural simple projection only, so possible Spin centers and represented central kernels do not create an exception. The theorem follows. \(\square\)

\subsection{Source and remaining-scope record}\label{source-and-remaining-scope-record}

The new Clifford construction uses Wonenburger's published Part II pp.48--49; its tensor-ideal proof is reproduced above. Its odd bimodule extension, norm preservation, and represented-quotient compatibility are proved here. The full Clifford-algebra extension theorem is not used. Finite Spin perfectness/quasisimplicity and its double-cover sequence are explicit published inputs at LST pp.1918--1919. The natural full-automorphism statement is the already completely reviewed D-OL07, including its precise Li--Xia source status. The character bound is the published LST Theorem 4.3.6 with Remark 4.3.7, already read in full.

No central quotient is presumed to be fixed by every natural automorphism. Only an automorphism actually compatible with the represented quotient is used, and (SP.10) proves its required kernel preservation. No Schur-multiplier classification is invoked. The theorem removes the genuine-Spin boundary only for odd \(q\), even \(d\ge1136\). Odd natural dimension and characteristic two remain separate.
